\documentclass[10pt,a4paper]{amsart}
\usepackage[margin=19mm]{geometry}
\usepackage{amsmath,amssymb,mathtools}
\usepackage[scr=rsfs,cal=euler]{mathalfa}
\usepackage{xfrac}
\usepackage[unicode,hidelinks,bookmarks=false]{hyperref}
\newcommand{\ie}{i.e.\ }
\newcommand{\cf}{cf.\ }
\newcommand{\resp}{respectively\ }
\newcommand{\Subsection}{\subsection}
\providecommand{\nobreakdash}{\nobreak\hbox{-}}
\setlength{\emergencystretch}{2em}
\multlinegap0pt
\usepackage{dsfont}
\usepackage{amssymb}
\usepackage{xcolor}
\usepackage{mathtools}
\newtheorem{prop}{Proposition}
\newtheorem{lem}{Lemma}
\newcommand{\R}{\mathbb{R}} 
\title{\small Sharp Weighted Discrete $p$-Hardy Inequality and Stability}
\author{Ali Barki}
\date{Source: arXiv:2501.00299v2 (CC BY 4.0)}
\begin{document}
\maketitle

\noindent\emph{Source and licence.} \small Sharp Weighted Discrete $p$-Hardy Inequality and Stability. Version arXiv:2501.00299v2 (CC BY 4.0), distributed by arXiv under the Creative Commons Attribution 4.0 International licence, http://creativecommons.org/licenses/by/4.0/, as recorded in the arXiv licence metadata for this exact version and shown on the abstract page https://arxiv.org/abs/2501.00299v2. Full licence notice: \texttt{licenses/arxiv-2501.00299-CC-BY-4.0.txt}.\par\smallskip

\noindent\textit{Editorial scope.} Counted unit: the article's prop 4.1 (source0.tex lines 656--663). The article's complete original proof of it is delivered with it (source0.tex lines 665--728, one \texttt{proof} environment of 64 lines). No mathematics is written, repaired or re-derived: the statement and the proof are the article's own lines, in the article's own notation, and no retained line is substituted or re-flowed.  No further source-local context is needed: every cross-reference of every retained line resolves inside the two delivered spans.  Nothing was omitted from any retained span, so the package carries no omission record.  No retained line contains a citation, so the wrapper carries no bibliography and no bibliography entry.  No retained line contains a figure environment, an image-inclusion command or drawing code, so no image is delivered and the package declares no asset dependency.  Editorial formatting only, in addition: an AMS \texttt{amsart} preamble that restates the article's own declarations and macros used by the retained lines, namely the environments \texttt{prop}, \texttt{lem} on separate counters, together with the standard packages those lines need.  The retained environments print in this document as Proposition 1, Lemma 1 in order of appearance, whereas the article prints the same items with the numbering of its own full document.  The complete native source of the article as distributed in the arXiv e-print for this exact version is bundled unchanged as \texttt{sources/170981/source0.tex}.
\begin{prop} \label{prop 4.1} For all $\gamma \in \mathbb{R}$, Let us consider, $\displaystyle S_n = \sum_{k=1}^n k^{\gamma - 1},
$ then, for all $n \geq 2$:
\begin{enumerate}
    \item[\textit{ i)}] For all  $1 < \gamma < 2$: $S_{n-1} + S_n < \frac{2 \, n^{\gamma}}{\gamma}$,
    \item[\textit{ ii)}] For all $\gamma > 0$: $S_{n-1} S_n < \frac{n^{2 \gamma}}{\gamma^2}$,
    \item[\textit{iii)}] For all  $\gamma \in \mathbb{R}$, we have $S_n > \exp\left(\frac{\gamma}{n}\right) S_{n-1}$.
\end{enumerate}
\end{prop}

\begin{proof} 
\textit{ i)}: since for $1 < \gamma < 2$ the function $t \to t^{\gamma-1}$ is concave on $(0,\infty)$ it follows that:
$$ S_{n-1}+S_{n} = 1+ \sum_{k=1}^{n-1} \left[ k^{\gamma-1} + (k+1)^{\gamma-1} \right]
\leq 1+ \sum_{k=1}^{n-1} 2 \int_{k}^{k+1} t^{\gamma-1} dt = 1+ 2 \int_{1}^{n} t^{\gamma-1} dt = \frac{\gamma-2}{\gamma} + \frac{2 n^{\gamma}}{\gamma} <  \frac{2 n^{\gamma}}{\gamma}  $$
\textit{ ii)}: starting with the case $1 < \gamma < 2,$ and using AM-GM inequality in addition to \textit{i)} gives us:
$$S_{n-1} S_{n} \leq  \left( \frac{S_{n-1}+S_{n}}{2} \right )^2  <  \frac{n^{2 \gamma}}{\gamma^2}$$
Now for $\gamma \in [0,1] \cup [2,\infty[,$ the function $t \to t^{\gamma-1}$ is convex therefore using Hermite-Hadamard inequality it follows that for all $k \in [|1,n|]:$
$$
k^{\gamma-1} = \left( \frac{k-\frac{1}{2}+  k+\frac{1}{2}}{2} \right)^{\gamma-1} \leq  \int_{k- \frac{1}{2}}^{k+ \frac{1}{2}} t^{\gamma-1} dt $$
By summing up we obtain that:
$$ S_{n} \leq \int_{\frac{1}{2}}^{n+\frac{1}{2}} t^{\gamma-1} dt \leq   \int_{0}^{n+\frac{1}{2}} t^{\gamma-1} dt  = \frac{(n+\frac{1}{2})^{\gamma}}{\gamma} $$
Thus, we may conclude that:
$ \displaystyle S_{n-1} S_{n} \leq  \frac{(n+\frac{1}{2} )^{\gamma} (n-\frac{1}{2})^{\gamma}}{\gamma^2} = \frac{(n^2-\frac{1}{4})^{\gamma}}{\gamma^2} < \frac{n^{2 \gamma}}{\gamma^2}$

\phantom{.} \\
\textit{(iii):}
 To address the last point, we analyze the cases based on the values of $\gamma:$ \\  \\
\underline{if $\gamma \leq 0,$} this is a trivial case since the inequality follows directly from \( S_{n} > S_{n-1} \). \\ \\ 
\underline{if $\gamma \in (0,1]\cup (1,2)$}: We start by establishing the following result: for all $x>0:$
$$ \frac{e^{x}}{e^{x}-1} - \frac{1}{2} - \frac{1}{x} > 0 $$
\phantom{.} \\
Define $\displaystyle f:x \to \frac{e^x}{e^x-1} - \frac{1}{2} - \frac{1}{x}.$ By analyzing its derivative, we obtain for all $x\in (0,\infty):$
$$ f'(x)= \frac{e^{x} (e^{x}+e^{-x}-x^2-2)}{x^2 (e^{x} -1)^2  }, $$ 
\phantom{.} \\ 
therefore, $f'(x)>0$ since $e^{x} + e^{-x} - x^2 - 2 = \displaystyle \sum_{k=2}^{\infty} \frac{x^{2k}}{(2k)!}> 0,$ \\ \\
hence for all $x\in (0,\infty):$ \ 
$\displaystyle f(x) > \lim\limits_{t\to 0^{+}} f(t) = \lim\limits_{t\to 0^{+}} \frac{t}{12} + o(t) = 0.$ Now for $\gamma \in (0,1],$ since $t\to t^{\gamma-1}$  is decreasing, we derive:
$$ \frac{S_n}{n^{\gamma-1}} = \frac{1}{n^{\gamma-1}}  \sum_{k=1}^{n} k^{\gamma-1} \leq  \frac{1}{n^{\gamma-1}}  \sum_{k=1}^{n} \int_{k-1}^{k} t^{\gamma-1} \ dt =  \frac{1}{n^{\gamma-1}} \int_{0}^{n} t^{\gamma-1} \ dt = \frac{n}{\gamma},
$$
Thus, we conclude that:
$$
\frac{S_n}{n^{\gamma-1}} \leq  \frac{n}{\gamma} < \left(\frac{\exp(\frac{\gamma}{n})}{\exp(\frac{\gamma}{n})-1}\right) -\frac{1}{2} < \frac{\exp(\frac{\gamma}{n})}{\exp(\frac{\gamma}{n})-1}.$$
which is equivalent to
$$ S_{n} > \exp(\frac{\gamma}{n})(S_{n}-n^{\gamma-1}) = \exp(\frac{\gamma}{n}) S_{n-1}.$$
For the case $\gamma \in (1,2)$ applying the result of \textit{ i)} along with the above inequality, we obtain:
$$ \frac{S_n}{n^{\gamma-1}} = \frac{1}{2 n^{\gamma-1}} (S_{n}+S_{n-1}) +\frac{1}{2} < \frac{n}{\gamma} + \frac{1}{2} < \frac{\exp(\frac{\gamma}{n})}{\exp(\frac{\gamma}{n})-1}.
$$
which is also equivalent to $S_{n} > \exp(\frac{\gamma}{n}) S_{n-1}.$ \\ \\
\underline{if $\gamma \geq 2$}: 
We will show by induction on \( n \) that for all $n \geq 1:$
$$ \frac{S_n}{n^{\gamma-1}} < \frac{\exp(\frac{\gamma}{n})}{\exp(\frac{\gamma}{n})-1}
$$
To complete the proof, we rely on the following lemma, whose proof is deferred to the end of the section.
\begin{lem} \label{lemma 4.2} Consider the function 
$F : [0,1]\times[2,\infty[ \longrightarrow  \R$ defined as follows:
$$\begin{array}{lrcl}
     &(x,\gamma) & \longmapsto & (1+x)^{\gamma-1} (1-exp(-\gamma x) ) - \exp \left( \frac{\gamma x}{1+x} \right) +1 \end{array}
$$
then for all $n \geq 1$ and $\gamma \geq 2,$ we have $F(\frac{1}{n} ,\gamma) \geq 0.$
\end{lem}
\phantom{.} \\
We proceed to prove the statement using induction.
For \( n = 1 \), the inequality holds trivially. Assuming that the inequality holds for \( n \), we aim to establish it for \( n + 1 \). First, by the induction hypothesis, we have:
$$ \frac{S_{n+1}} {(n+1)^{\gamma-1}} = ( 1+ \frac{1}{n})^{1-\gamma} \,  \frac{S_{n}}{n^{\gamma-1}} +1 <
( 1+ \frac{1}{n})^{1-\gamma} \, \left[ \frac{\exp(\frac{\gamma}{n})}{\exp(\frac{\gamma}{n})-1} \right] + 1$$
From Lemma \ref{lemma 4.2}, we have the following equivalent formulation
$$( 1+ \frac{1}{n})^{1-\gamma}  \, \left[ \frac{\exp(\frac{\gamma}{n})}{\exp(\frac{\gamma}{n})-1} \right]  \leq \frac{1}{\exp(\frac{\gamma}{n+1})-1}$$
By combining the induction hypothesis with the result of Lemma \ref{lemma 4.2}, we derive:
$$\frac{S_{n+1}} {(n+1)^{\gamma-1}}
<   \frac{1}{\exp(\frac{\gamma}{n+1})-1} +1  =  \frac{\exp(\frac{\gamma}{n+1})}{\exp(\frac{\gamma}{n+1})-1}.$$ Hence, the induction step is verified, and the inequality holds for all $n \geq 1:$ (with the convention $S_{0}=0$).
$$ \frac{S_n}{n^{\gamma-1}} < \frac{\exp(\frac{\gamma}{n})}{\exp(\frac{\gamma}{n})-1} \Leftrightarrow S_{n} > \exp(\frac{\gamma}{n}) S_{n-1}.
$$
This concludes the proof of part \textit{iii)} and completes the proposition.
\end{proof}

\end{document}
